\documentclass[11pt,a4paper]{article}
\usepackage[margin=25mm]{geometry}
\usepackage{fontspec}
\setmainfont{TeX Gyre Pagella}
\usepackage{amsmath,amssymb,mathtools}
\usepackage{graphicx}
\usepackage{xcolor}
\usepackage[unicode,colorlinks=true,linkcolor=blue,urlcolor=blue]{hyperref}
\usepackage{bookmark}
\usepackage{enumitem}
\setlist{itemsep=3pt,topsep=5pt}
\setlength{\emergencystretch}{3em}
\newcommand{\tightlist}{\setlength{\itemsep}{0pt}\setlength{\parskip}{0pt}}
\newcommand{\passthrough}[1]{#1}
\newcommand{\pandocbounded}[1]{#1}
\makeatletter
\def\maxwidth{\ifdim\Gin@nat@width>\linewidth\linewidth\else\Gin@nat@width\fi}
\def\maxheight{\ifdim\Gin@nat@height>0.75\textheight0.75\textheight\else\Gin@nat@height\fi}
\makeatother
\setkeys{Gin}{width=\maxwidth,height=\maxheight,keepaspectratio}
\hypersetup{pdfauthor={OpenAI GPT-6.1 Sol, Codex, Ultra effort}}
\begin{document}
\section{The Chebotarev density
theorem}\label{the-chebotarev-density-theorem}

\emph{Written by OpenAI GPT-6.1 Sol in Codex, Ultra effort, October
2026. Self-checked by the writing AI; no independent review is claimed.
Public domain (CC0).}

\textbf{Lesson 22.} Every conjugacy class of a finite Galois group
occurs at infinitely many unramified primes. More precisely, its share
of the primes, measured by Dirichlet density, is its share of the group.
We prove this first for abelian groups by Hecke characters, then count
degree-one primes in a cyclic fixed-field extension to obtain the
general theorem. Its ray-class corollary supplies the prime needed for
the full quadratic-form induction in section 6A, starting from the conic
theorem of lesson 15.

The arithmetic Frobenius and Artin L-function conventions are those of
\href{../artin-l-functions-conductors-and-discriminants.html}{Artin
L-functions, conductors and discriminants}. The analytic number-field
inputs are the written Theorem 10.1 of
\href{https://kokunoyumeto.github.io/open-math-courses/courses/NT-CFT/proof-dependencies.html\#NT-ADL-10}{Hecke
L-functions and the Dedekind zeta function}, which makes a nontrivial
finite-order Hecke L-function holomorphic at 1, and Theorem 16.2 of
\href{https://kokunoyumeto.github.io/open-math-courses/courses/NT-CFT/proof-dependencies.html\#NT-ANT-16}{The
Dedekind zeta function and the analytic class number formula}, which
gives the positive simple zeta pole there. Neither input uses
Chebotarev. The function-field versions needed here follow directly from
section 6 of the preceding lesson, as checked below. We use the written
existence and ray correspondence of lessons 17--18.

\textbf{Prerequisite proof availability.} The named results below
identify specific programme lessons. The
\href{https://kokunoyumeto.github.io/open-math-courses/courses/NT-CFT/proof-dependencies.html\#lesson-22}{prerequisite
record} distinguishes published proofs, supplied owner texts awaiting
publication, and missing full proofs. A record or external reference is
not a supplied proof; arguments using an unavailable prerequisite retain
that dependency.

\subsection{1. Which density is being
computed?}\label{which-density-is-being-computed}

For a global field \(K\), let \(\mathcal P_K\) denote its finite places
and write \(N\mathfrak p\) for the residue cardinality. For
\(S\subseteq\mathcal P_K\), define its \textbf{Dirichlet density}, when
the limit exists, by \[
\delta_K(S)=\lim_{s\downarrow1}
\frac{\sum_{\mathfrak p\in S}(N\mathfrak p)^{-s}}
     {\log(1/(s-1))}.
\tag{1}
\] Here \(s\) is real and tends to 1 from above. Its \textbf{natural
density}, when it exists, is \[
\lim_{X\to\infty}
\frac{\#\{\mathfrak p\in S:N\mathfrak p\leq X\}}
     {\#\{\mathfrak p\in\mathcal P_K:N\mathfrak p\leq X\}}.
\tag{2}
\] The two are different definitions. Our Chebotarev theorem asserts
(1).

\textbf{Proposition 22.1.} The full set of primes has Dirichlet density
1. Densities lie in \([0,1]\), are monotone, are unchanged by finite
modifications, and are finitely additive on disjoint sets whose
densities exist. Complements have density \(1-\delta\). More generally a
set whose prime sum remains bounded as \(s\downarrow1\) has density
zero. Whenever the natural density (2) exists, its Dirichlet density
exists and equals it.

\textbf{Proof.} The zeta Euler product gives \[
\log\zeta_K(s)=\sum_{\mathfrak p}(N\mathfrak p)^{-s}
 +\sum_{\mathfrak p}\sum_{m\geq2}\frac{(N\mathfrak p)^{-ms}}m.
\tag{3}
\] For \(s\geq1\), the last sum is at most
\(2\sum_{\mathfrak p}(N\mathfrak p)^{-2}\), a convergent subseries of
\(\zeta_K(2)\). The positive simple pole at 1 gives \[
\sum_{\mathfrak p}(N\mathfrak p)^{-s}
=\log\frac1{s-1}+O_K(1).
\tag{4}
\] For number fields the zeta input is the exact written theorem
identified above. For function fields, use the test
\(f=\prod_v\mathbf1_{\mathcal O_v}\) and the trivial character in
(29)--(30) of the preceding lesson. Its zeta integral is \(\zeta_K(s)\).
The positive-level parts are finite Laurent polynomials. At \(s=1\) the
first correction has a simple pole with residue
\(\kappa\widehat f(0)/\log R>0\); the second correction is regular.
Indeed \(\widehat f(0)=\operatorname{vol}(\prod_v\mathcal O_v)>0\), with
only finitely many local volumes different from 1. Thus the pole is
simple and positive, proving (4) here too.

Replacing the denominator of (1) by the full prime sum consequently
gives the same limit. All assertions about bounds, inclusion, disjoint
sums, complements and bounded contributions now follow from nonnegative
summands and elementary limit laws. Finite contributions are bounded.
Countable additivity is not asserted: a countable union of individual
zero-density primes is the full prime set.

For the last assertion, write \(A_S(t),A_K(t)\) for the two prime counts
and suppose \(A_S(t)/A_K(t)\to d\). Partial summation gives \[
\sum_{\mathfrak p\in S}(N\mathfrak p)^{-s}
=s\int_1^\infty A_S(t)t^{-s-1}\,dt
\qquad(s>1).
\tag{5}
\] The endpoint at infinity vanishes: choose \(1<\sigma<s\); convergence
of the prime sum at \(\sigma\) bounds \(A_K(t)=O(t^\sigma)\). For any
\(\epsilon>0\), choose \(T\) beyond which
\(|A_S(t)-dA_K(t)|\leq\epsilon A_K(t)\). The integral of the difference
on \([1,T]\) stays bounded as \(s\downarrow1\); the tail is at most
\(\epsilon\) times the full prime sum. Divide by that diverging sum in
(4) and let \(\epsilon\to0\). This proves the equality of densities.
\(\square\)

A useful relative version of a zero-density observation will be used in
the proof. If \(E/K\) is finite and \(u\) is a prime of \(E\) of residue
degree \(f(u/v)>1\) over \(K\), then \(Nu=(Nv)^{f(u/v)}\). There are at
most \([E:K]\) primes over \(v\), and therefore \[
\sum_{\substack{u\in\mathcal P_E\\f(u/v)>1}}(Nu)^{-s}
\leq[E:K]\sum_{v\in\mathcal P_K}(Nv)^{-2s}=O_{E/K}(1)
\quad(s\downarrow1).
\tag{6}
\] Thus these relative degree-greater-than-one primes contribute no
Dirichlet density. In particular this argument does not discard
degree-one primes merely because their degree over the prime field is
larger than one.

\subsection{2. Finite-order Hecke characters do not vanish at
1}\label{finite-order-hecke-characters-do-not-vanish-at-1}

\textbf{Proposition 22.2.} If \(\omega:C_K\to\mathbf C^\times\) is a
nontrivial finite-order continuous Hecke character, then its primitive
L-function is holomorphic and nonzero at 1.

\textbf{Proof.} The kernel is open of finite index. Global existence
gives its finite abelian class field \(L/K\), and reciprocity identifies
\(\omega\) with a nontrivial character of \(G=\operatorname{Gal}(L/K)\).
For every nontrivial character \(\chi\) of \(G\), Theorem 21.2
identifies its Artin L-function with its primitive finite-order Hecke
L-function. For number fields the Hecke theorem makes it holomorphic at
1: a nontrivial finite-order character cannot be a pure imaginary norm
twist, since the norm direction is connected and its finite image must
be trivial.

For function fields the same holomorphy follows from the explicit
corrections in the preceding lesson. If \(\chi\) is nontrivial on
\(C_K^1\), its correction is zero and its zeta integral is a finite
Laurent polynomial. If it is trivial there, it has
\(\chi(c)=\lambda\ne1\), because \(c\) generates the norm quotient and
the global character is nontrivial. At \(s=1\), both denominators
\(1-\lambda^{-1}\) and \(1-\lambda^{-1}R^{-1}\) are nonzero. Thus its
integral and primitive L-function are holomorphic there. This checks
explicitly the possible constant-field characters.

Because \(G\) is abelian, all irreducible characters have dimension one.
Proposition 21.1 gives \[
\zeta_L(s)=\zeta_K(s)\prod_{\substack{\chi\in\widehat G\\\chi\ne1}}L_K(s,\chi).
\tag{7}
\] Both zeta functions have simple poles with positive residues at 1.
Their quotient is holomorphic and has nonzero value equal to the
quotient of those residues. Every factor on the right of the quotient
identity is holomorphic at 1, so none can vanish there. In particular
\(L_K(1,\omega)\ne0\). This uses no density theorem or assumption about
primes in a ray class. \(\square\)

If extra primes are omitted to describe an imprimitive character modulo
a larger ray modulus, its L-function is multiplied by finitely many
factors \(1-\chi(\operatorname{Frob}_v)(Nv)^{-s}\). At 1 these are
nonzero, because \(|\chi(\operatorname{Frob}_v)|=1<Nv\). The same
nonvanishing statement holds for that imprimitive product, but the
primitive identity (7) keeps its Euler factors explicit.

\subsection{3. The abelian density
theorem}\label{the-abelian-density-theorem}

\textbf{Theorem 22.3.} Let \(L/K\) be a finite abelian extension with
group \(G\). For each \(\sigma\in G\), the unramified primes with
arithmetic Frobenius \(\sigma\) have Dirichlet density \(1/|G|\).

\textbf{Proof.} Remove the finitely many ramified primes, and for
\(\chi\in\widehat G\) put \[
S_\chi(s)=\sum_{v\text{ unramified}}\chi(\operatorname{Frob}_v)(Nv)^{-s}.
\tag{8}
\] The logarithmic Euler expansion differs from this first-power sum by
a uniformly bounded sum near 1: the terms of power at least two are
bounded exactly as in (3), and the removed finite places contribute
bounded logarithmic factors. For \(\chi\ne1\), Proposition 22.2 permits
a holomorphic logarithm of \(L_K(s,\chi)\) in a disk about 1. On the
real segment \(1<s<1+\epsilon\), this logarithm differs from the
absolutely convergent Euler logarithm by a constant multiple of
\(2\pi i\), since both are continuous logarithms of the same nonzero
function. It follows that \[
S_\chi(s)=O(1)\quad(\chi\ne1),\qquad
S_1(s)=\log\frac1{s-1}+O(1).
\tag{9}
\] Finite abelian character orthogonality gives \[
\mathbf1_{g=\sigma}=\frac1{|G|}\sum_{\chi\in\widehat G}
\overline{\chi(\sigma)}\chi(g).
\tag{10}
\] For clarity, if \(g\sigma^{-1}\ne1\), choose a character nontrivial
on it; multiplying the character sum by that value both permutes its
terms and changes its value, forcing it to be zero. If \(g=\sigma\), all
terms are 1. Characters separate elements by the written finite duality
theorem used in lessons 3 and 18. Multiply (10) by \((Nv)^{-s}\), sum
over primes, and insert (9). Only the trivial character contributes the
logarithmic main term. Hence the desired prime sum is
\(|G|^{-1}\log(1/(s-1))+O(1)\), which proves the density. \(\square\)

\subsection{4. The degree-one count in a cyclic fixed
field}\label{the-degree-one-count-in-a-cyclic-fixed-field}

\textbf{Theorem 22.4 (Chebotarev).} Let \(L/K\) be finite Galois, with
group \(G\), and let \(C\) be a conjugacy class in \(G\). The unramified
primes of \(K\) whose arithmetic Frobenius conjugacy class is \(C\) have
Dirichlet density \[
\delta_K(\mathcal P_C)=\frac{|C|}{|G|}.
\tag{11}
\]

\textbf{Proof.} Choose \(\sigma\in C\), put \(H=\langle\sigma\rangle\),
\(M=L^H\), and \(Z=Z_G(\sigma)\). The extension \(L/M\) is cyclic with
group \(H\). Theorem 22.3 says that the primes \(u\) of \(M\) with
Frobenius exactly \(\sigma\) have density \(1/|H|\). Omit the primes
over the finite ramification set of \(L/K\), and also the relative
degree-greater-than-one primes. Formula (6) shows that the latter
omission changes the weighted prime sum by \(O(1)\). For the remaining
primes, \(Nu=Nv\) for \(v=u\cap K\).

We count how many such \(u\) lie over a given unramified \(v\). Choose a
Frobenius \(\tau\in G\) at \(v\). Its action on the embeddings of \(M\),
identified with \(G/H\), has cycles of lengths equal to the residue
degrees of primes over \(v\). One can see this directly by completing:
the orbits of the decomposition group \(\langle\tau\rangle\) correspond
to the completions of \(M\), and the orbit length is their unramified
residue degree. A degree-one prime is a fixed coset \(gH\), equivalent
to \(g^{-1}\tau g\in H\). Its Frobenius in \(L/M\) is this element of
\(H\); it is independent of the coset representative because \(H\) is
cyclic.

Consequently a degree-one prime contributes to the cyclic set with
Frobenius \(\sigma\) precisely when \[
g^{-1}\tau g=\sigma.
\tag{12}
\] If \(\tau\notin C\), there are none. If \(\tau\in C\), the set of
solutions \(g\) has cardinality \(|Z|\): after choosing one solution
\(g_0\), all others are \(g_0z\), \(z\in Z\). Each right coset of \(H\)
accounts for exactly \(|H|\) of these solutions, since \(H\subseteq Z\).
Thus there are exactly \(|Z|/|H|\) contributing primes of \(M\) over
every \(v\in\mathcal P_C\).

Taking weighted sums gives \[
\sum_{\substack{u\in\mathcal P_M\\\operatorname{Frob}_u(L/M)=\sigma}}(Nu)^{-s}
=\frac{|Z|}{|H|}\sum_{v\in\mathcal P_C}(Nv)^{-s}+O(1).
\tag{13}
\] Divide by \(\log(1/(s-1))\). The left side has limit \(1/|H|\), so
the sum on the right has limit \(1/|Z|\). The conjugation orbit formula
says \(|C|=|G|/|Z|\), proving (11). \(\square\)

This is Deuring's cyclic fixed-field reduction. Counting all primes of
\(M\) as if they had norm \(Nv\) would lose the residue-degree condition
and invalidate (13). Formula (6) is exactly what permits their removal.
The argument works over function fields as well as number fields;
constant extensions may restrict the allowed degrees, but the stated
Dirichlet density remains (11).

For a union of conjugacy classes, finite additivity gives the density
equal to the size of the union divided by \(|G|\). In particular every
nonempty such set of primes is infinite: a finite set has density zero,
whereas the computed density is positive.

\subsection{5. Fields and ray classes determined by
primes}\label{fields-and-ray-classes-determined-by-primes}

\textbf{Corollary 22.5 (Bauer).} Two finite Galois extensions of \(K\)
inside a common separable closure are equal if their sets of completely
split primes agree outside a finite set. More generally, if \(E/K\) is
Galois and \(M/K\) is any finite separable extension, then \[
E\subseteq M
\quad\Longleftrightarrow\quad
\text{almost every prime admitting a degree-one prime in }M
\text{ splits completely in }E.
\tag{14}
\]

\textbf{Proof.} Take a finite Galois \(N/K\) containing \(E,M\), with
group \(G\), and put \(H_E=\operatorname{Gal}(N/E)\),
\(H_M=\operatorname{Gal}(N/M)\). The first is normal. Outside the finite
ramification set, a prime splits completely in \(E\) exactly when its
Frobenius lies in \(H_E\). It admits a degree-one prime in \(M\) exactly
when the Frobenius fixes a coset of \(G/H_M\), equivalently when its
conjugacy class meets \(H_M\).

If \(E\subseteq M\), then \(H_M\subseteq H_E\); every class meeting
\(H_M\) lies entirely in \(H_E\), proving the forward implication.
Conversely, if some \(\tau\in H_M\setminus H_E\) exists, its class meets
\(H_M\) and is disjoint from \(H_E\), by normality. Chebotarev supplies
infinitely many primes of this class, contradicting the asserted
almost-everywhere inclusion. Thus \(H_M\subseteq H_E\), so
\(E\subseteq M\). Applying this with the two Galois fields in both
directions proves the equality assertion. \(\square\)

\textbf{Corollary 22.6 (primes in ray classes).} For a number field
\(K\) and a modulus \(\mathfrak m\), every ideal ray class contains
infinitely many primes outside any prescribed finite set. More precisely
its prime set has Dirichlet density
\(1/|\operatorname{Cl}_{\mathfrak m}(K)|\). Every element of the group
of any finite abelian extension of a global field similarly occurs at
infinitely many primes outside a prescribed finite set.

\textbf{Proof.} The written ray-class theorem identifies
\(\operatorname{Cl}_{\mathfrak m}(K)\) with the Galois group of its ray
field and sends a prime away from \(\mathfrak m\) to its arithmetic
Frobenius. Theorem 22.3 computes its density. Adding any finite
exceptional set changes no density by Proposition 22.1, so infinitely
many remain. The final assertion is Theorem 22.3 applied to the
specified finite abelian extension, with the same finite-omission
argument. \(\square\)

For \(\mathbf Q\) with modulus \(m\infty\), the ray group is
\((\mathbf Z/m\mathbf Z)^\times\), and a prime's Artin image is its
residue class. Thus every progression \(a\bmod m\), \((a,m)=1\), has
prime Dirichlet density \(1/\varphi(m)\) and contains infinitely many
primes. This is Dirichlet's theorem, with the density specified.

Corollary 22.6 lets us choose a prime outside a bad set with Frobenius
equal to the inverse of a prescribed idèle class. Section 6A uses this
to prove Hasse--Minkowski. The quadratic-form argument does not enter
the analytic or character proofs above.

\subsection{6. Cubic examples}\label{cubic-examples}

For \(L=\mathbf Q(\sqrt[3]2,\zeta_3)\), the preceding lesson proved
\(G=S_3\). Outside 2 and 3, reduction of \(X^3-2\) has distinct roots in
the splitting residue field. Arithmetic Frobenius permutes them, and the
length of each orbit equals the degree of the corresponding irreducible
factor over \(\mathbf F_p\): a root of degree \(d\) is fixed by the
\(d\)-th Frobenius power and by no smaller positive power. Thus
Chebotarev gives:

{\def\LTcaptype{none} % do not increment counter
\begin{longtable}[]{@{}
  >{\raggedright\arraybackslash}p{(\linewidth - 6\tabcolsep) * \real{0.2143}}
  >{\raggedleft\arraybackslash}p{(\linewidth - 6\tabcolsep) * \real{0.2857}}
  >{\raggedright\arraybackslash}p{(\linewidth - 6\tabcolsep) * \real{0.2143}}
  >{\raggedleft\arraybackslash}p{(\linewidth - 6\tabcolsep) * \real{0.2857}}@{}}
\toprule\noalign{}
\begin{minipage}[b]{\linewidth}\raggedright
Frobenius class
\end{minipage} & \begin{minipage}[b]{\linewidth}\raggedleft
Number of elements
\end{minipage} & \begin{minipage}[b]{\linewidth}\raggedright
Factor degrees of \(X^3-2\bmod p\)
\end{minipage} & \begin{minipage}[b]{\linewidth}\raggedleft
Dirichlet density
\end{minipage} \\
\midrule\noalign{}
\endhead
\bottomrule\noalign{}
\endlastfoot
Identity & 1 & \(1+1+1\) & \(1/6\) \\
Transpositions & 3 & \(1+2\) & \(1/2\) \\
Three-cycles & 2 & \(3\) & \(1/3\) \\
\end{longtable}
}

In the first row the prime splits completely in the splitting field. In
the third row the cubic is irreducible modulo \(p\). The finite
exceptional primes affect none of these densities.

\subsection{6A. General quadratic
forms}\label{a.-general-quadratic-forms}

\textbf{Theorem 22.7 (Hasse--Minkowski).} A nondegenerate quadratic form
over a number field is isotropic over that field if and only if it is
isotropic over every completion.

\textbf{Proof.} The inputs are now available: Hasse's norm theorem and
the conic principle in lesson 15; global reciprocity with local
compatibility in lesson 16; existence in lesson 17; ray fields in lesson
18; and Corollary 22.6 above. The latter lets us prescribe a Frobenius
element outside any finite set. Their proofs do not use the
quadratic-form theorem.

Dimension \(1\) is immediate. In dimension \(2\), local isotropy says a
fixed element is a square everywhere. Its quadratic extension would
split completely at every prime, contradicting Corollary 14.4 unless the
extension is trivial. Dimension \(3\) is Corollary 15.4 of
\href{../the-norm-index-bound-and-hasses-norm-theorem.html}{The norm
index bound and Hasse's norm theorem}.

Suppose \(n\geq4\), diagonalize, and write \(q=q_1+q_2\), with \(q_1\)
binary and \(q_2\) of dimension \(n-2\geq2\). At every place, choose a
nonzero \(t_v\) represented by both \(q_1\) and \(-q_2\). Local isotropy
supplies such a value unless its two part-values are zero. In that case
at least one part is isotropic and contains a hyperbolic plane, hence
represents every scalar; choose a nonzero value of the other part.

The simultaneous value set contains a neighborhood of \(t_v\). At a
representation of a nonzero value, some coordinate \(z\ne0\). Changing
just that coordinate changes the value by \(2czh+ch^2\), which maps a
sufficiently small ball onto a neighborhood by contraction at finite
places, and the usual one-variable inverse function argument at
archimedean places. At a real place all multiples by positive scalars
are represented, by scaling coordinates; at a complex place every
nonzero scalar is represented.

Choose \(S\) containing the archimedean and dyadic places, all nonunit
diagonal coefficients, and the ramified places of the binary norm
fields. If \(q_1=aX^2+bY^2\), its field is \(M_1=K(\sqrt{-b/a})\),
omitting it if already split. If \(q_2\) is binary, also take the field
for \(-q_2\), with leading coefficient \(a_2\). Outside \(S\), a binary
unit form represents every unit: over the odd residue field it is either
split or a scalar times the quadratic norm, which is surjective; a
nonsingular residue representation lifts. A unit form of dimension at
least \(3\) is isotropic locally. To see this, in three variables the
two sets of values \(c_1x^2\) and \(-c_3-c_2y^2\), each of size
\((q_v+1)/2\), intersect; the solution with third coordinate \(1\)
lifts. An isotropic nondegenerate form contains a hyperbolic plane and
represents all scalars.

Let \(x\) be the idèle with components \(t_v\) in \(S\) and \(1\)
elsewhere. For each binary norm field \(M_i\), representation says
\(t_v/a_i\) is a local norm at \(S\). Global reciprocity for the
principal element \(a_i\), whose other local symbols are trivial since
it is a unit at unramified places, gives \[
\operatorname{rec}_{M_i/K}(x)
=\prod_{v\in S}\operatorname{rec}_{M_i/K,v}(t_v)
=\prod_{v\in S}\operatorname{rec}_{M_i/K,v}(a_i)=1.
\tag{34}
\] Choose a product subgroup \(U\subset J_K\), consisting of
sufficiently deep principal units at the finite places in \(S\),
positive reals at real places, full complex groups, and units outside
\(S\). Require that \(t_vU_v\) lies in both local value sets and \(U_v\)
lies in each binary local norm group. Its image \(H\subset C_K\) is open
of finite index: it has full content image and its remaining quotient
comes from compact \(C_K^1\). The ray class correspondence gives its
abelian class field \(A/K\), containing the \(M_i\).

The ray-prime theorem chooses a prime \(\mathfrak p\notin S\) with \[
\operatorname{Frob}_{\mathfrak p}(A/K)
=\operatorname{rec}_{A/K}(x)^{-1}.
\] By (34) it splits in every \(M_i\). If \(\pi_{\mathfrak p}\) is its
idèle uniformizer, then \(x\pi_{\mathfrak p}=t u\) for some
\(t\in K^\times\) and \(u\in U\). Thus \(t\) is in the chosen
simultaneous value neighborhood at \(S\), is a unit outside
\(S\cup\{\mathfrak p\}\), and has valuation \(1\) at \(\mathfrak p\). At
this last prime every binary part is split, and every larger part is
isotropic. Consequently \(t\) is represented by \(q_1\) and by \(-q_2\)
at every place.

Apply induction to \(q_1-tZ^2\) and \(q_2+tZ^2\), of dimensions \(3\)
and \(n-1\). Each is globally isotropic. If its \(Z\)-coordinate is
nonzero, division produces the required global representation. If that
coordinate is zero, its original part is isotropic and represents every
scalar via a hyperbolic plane, so again it represents the required
value. Combining the representations gives a nonzero global isotropic
vector of \(q\). This completes the induction.

The hyperbolic-plane assertion used here is elementary: given an
isotropic \(e\ne0\), choose \(f\) with the symmetric bilinear form
\(B(x,y)=(q(x+y)-q(x)-q(y))/2\) and \(B(e,f)=1\), and replace \(f\) by
\(f-q(f)e/2\). Their plane has form \(2XY\). \(\square\)

\subsection{7. Exercises and complete
solutions}\label{exercises-and-complete-solutions}

\subsubsection{Exercise 1 --- Complete splitting in the cubic splitting
field
(easy)}\label{exercise-1-complete-splitting-in-the-cubic-splitting-field-easy}

Compute the density of primes splitting completely in
\(\mathbf Q(\sqrt[3]2,\zeta_3)\).

\textbf{Solution.} Eisenstein irreducibility and discriminant \(-108\),
a nonsquare, give splitting group \(S_3\) of order six. At an unramified
prime complete splitting means trivial decomposition group, equivalently
identity Frobenius. The identity class has one element, so Theorem 22.4
gives density \(1/6\). Removing the ramified primes 2 and 3 makes no
difference to the limit.

\subsubsection{Exercise 2 --- Irreducible cubic reductions
(medium)}\label{exercise-2-irreducible-cubic-reductions-medium}

Compute the density of rational primes at which \(X^3-2\) is irreducible
modulo \(p\).

\textbf{Solution.} At \(p\ne2,3\), the polynomial is separable modulo
\(p\). Its factor degrees are the lengths of Frobenius orbits on its
three roots, by the finite-field orbit argument in section 6. It is
irreducible exactly when Frobenius is a three-cycle. The two
three-cycles form one conjugacy class in \(S_3\); hence its density is
\(2/6=1/3\). At 2 the reduction is \(X^3\), and at 3 it is \((X+1)^3\),
so neither is irreducible, and their exclusion in any event contributes
zero density.

\subsubsection{Exercise 3 --- Recovering a Galois field from split
primes
(medium)}\label{exercise-3-recovering-a-galois-field-from-split-primes-medium}

Prove the equality assertion in Bauer's theorem directly from
Chebotarev.

\textbf{Solution.} Put the two Galois fields \(E,F\) in their Galois
compositum \(N\). Let \(H_E,H_F\) be the kernels of restriction from
\(G=\operatorname{Gal}(N/K)\). If the fields are distinct, the kernels
are distinct, so some element belongs to one and not the other. Both
kernels are normal, so its entire conjugacy class belongs to that same
difference. Chebotarev gives that class positive density and therefore
infinitely many unramified primes. Every one splits completely in the
first field and fails to split completely in the second. This
contradicts agreement of the splitting sets outside a finite set. Thus
the kernels and their fixed fields are equal. This argument also
explains why a mere finite list of matching split primes cannot
determine the field.

\subsubsection{Exercise 4 --- The pole argument without a hidden zero
(hard)}\label{exercise-4-the-pole-argument-without-a-hidden-zero-hard}

Prove Proposition 22.2 in detail, explaining why the zeta factorization
forbids a zero of any nontrivial character and why extra omitted Euler
factors do not cause an exception.

\textbf{Solution.} Let \(L/K\) be the class field of \(\ker\omega\). Its
group is finite abelian, and the primitive Artin--Hecke comparison
identifies every group character \(\chi\) with its finite-order Hecke
L-function. The analytic input gives holomorphy at 1 for every
nontrivial character. Over number fields this is the finite-order case
of the written Hecke theorem. Over function fields it follows from the
finite-polynomial and correction formulas: a character nontrivial on
norm-one classes has no correction, and a nontrivial norm character has
\(\lambda\ne1\), so its denominators are nonzero at 1. The trivial
character is precisely \(\zeta_K\), which has a simple pole with
positive residue \(c_K\). The same zeta argument for \(L\) gives a
positive simple pole with residue \(c_L\).

In a punctured neighborhood of 1, divide (7) by \(\zeta_K\). The
quotient extends holomorphically with value \(c_L/c_K\ne0\). Thus the
finite product of the nontrivial holomorphic character L-functions has a
nonzero value. If any one factor vanished, the product would vanish;
other factors cannot cancel it by poles because they have no poles at 1.
Equivalently, the total order at 1 in (7) is
\(-1=-1+\sum_{\chi\ne1}\operatorname{ord}_1L(s,\chi)\), with every
summand nonnegative, so all are zero. This proves nonvanishing for each
character, including \(\omega\).

A larger ray modulus deletes finitely many primitive factors,
multiplying by \(1-\chi(\operatorname{Frob}_v)(Nv)^{-s}\) at
character-unramified omitted primes. Each has absolute-value input
\((Nv)^{-1}<1\) at 1 and cannot be zero there. Character-ramified primes
already have primitive factor 1. Consequently neither holomorphy nor
nonvanishing at 1 is altered by these extra omissions. No prime-density
or ray-prime existence assertion was used in this argument.

\subsection{What this lesson does not
prove}\label{what-this-lesson-does-not-prove}

The number-field Hecke continuation and positive zeta-pole theorems have
the exact written internal providers identified at the beginning; the
necessary function-field consequences were checked explicitly from the
preceding lesson's proof. The global class-field and ray-field
constructions were proved in lessons 17--18. The present lesson proves
nonvanishing by their zeta factorization and supplies the character
sums, cyclic fixed-field count, densities and applications in full.

The theorem proved here is Dirichlet-density Chebotarev. We make no
assertion of natural density or an effective error term. Those require
additional prime-counting analysis. Proposition 22.1 proves only that an
existing natural density must agree with the Dirichlet density; it
supplies no converse.

\subsection{Editable edition}\label{editable-edition}

The reading edition provides the complete LaTeX source of this lesson,
the cumulative course LaTeX and the editable source ZIP. The archive
contains all twenty-four Markdown lessons, complete LaTeX bodies,
original diagrams, metadata and reproduction instructions.

\subsection{References}\label{references}

The lesson proves nonvanishing at 1, abelian character density and the
cyclic fixed-field coset count, then Bauer and primes in ray classes. It
proves Dirichlet density over both sorts of global field; no
natural-density or effective assertion is substituted.

\begin{itemize}
\tightlist
\item
  \href{https://www.jmilne.org/math/CourseNotes/CFT.pdf}{J. S. Milne,
  Class Field Theory, version 4.03}.
\item
  \href{https://math.mit.edu/~poonen/786/notes.pdf}{Bjorn Poonen, Tate's
  Thesis, MIT 18.786 lecture notes (2015)}.
\item
  \href{https://kskedlaya.org/cft/sec_abstractcft1.html}{Kiran S.
  Kedlaya, Notes on class field theory, author-hosted HTML edition}.
\end{itemize}

The \href{../free-proofs.html}{proof guide} gives the lesson sequence
and the exact prerequisite record. External references accompany the
written arguments.

\end{document}
